% Compilar com o XeLaTeX
\documentclass{gjvnq-livreto}
\usepackage{ccicons}

\title{Classe \LaTeX: gjvnq-livreto}
\author{G. Queiroz\\<\mailto{gabrieljvnq@gmail.com}>}
\date{2019}

\begin{document}
\maketitle
\vspace{10em}
\begin{center}
\BeginAccSupp{ActualText=CC-BY-SA}
\ccbysa
\EndAccSupp{}

\footnotesize
\link{Licença Creative Commons Atribuição-CompartilharIgual 4.0}{https://creativecommons.org/licenses/by-sa/4.0/}
\end{center}
\clearpage
\tableofcontents
\clearpage


$x = \frac{-b \pm \sqrt{b^2 - 4ac}}{a}$

Luís argüia à Júlia que «brações, fé, chá, óxido, pôr, zângão» eram palavras do português.

À noite, vovô Kowalsky vê o ímã cair no pé do pinguim queixoso e vovó põe açúcar no chá de tâmaras do jabuti feliz. (90 letras, incluindo todas as letras acentuadas)

% \begin{center}
% \includegraphics[width=1cm]{CC-BY-SA.ps}\\
% \footnotesize
% Documento sob a \link{Licença Creative Commons Atribuição-CompartilharIgual}{https://creativecommons.org/licenses/by-sa/4.0/}.
% \end{center}

\chapter{Básico}

\section{Letras e Acentos}

As letras do alfabeto da língua portuguesa, incluindo os acentos, funcionam como se espera:

\begin{latexExample}
Luís argüia à Júlia que «brações, fé, chá, óxido, pôr, zângão» eram palavras do português.
\end{latexExample}

As seguintes letras gregas também podem ser digitadas diretamente: α, β, γ, δ, ε, ζ, η, θ, ι, κ, λ, μ, ν, ξ, ο, π, ρ, σ, τ, υ, φ, χ, ψ, ω, Α, Β, Γ, Δ, Ε, Ζ, Η, Θ, Ι, Κ, Λ, Μ, Ν, Ξ, Ο, Π, Ρ, Σ, Τ, Υ, Φ, Χ, Ψ, Ω, ς, φ.

\subsection{Outro Símbolos de Texto}
\begin{codeBox}
\begin{longtable}{lll}
Real & \verb|\BRL| & \BRL \\
Euro & \verb|\EUR| & \EUR \\
Porcento & \verb|\%| & \% \\
Número & \verb|\textnumero| & \textnumero \\
Interrobang & \verb|\textinterrobang| & \textinterrobang \\
Interrobang invertido & \verb|\textinterrobangdown| & \textinterrobangdown \\
\end{longtable}
\end{codeBox}

\section{Matemática}

% \begin{latexExample}
% Prove que:
% \begin{equation}
% e^{\pi i} = -1
% \end{equation}
% \end{latexExample}

% \begin{latexExample}
% Prove que a soma dos $n$ primeiros termos de uma P.G. com razão $q$, pode ser expressa como:
% \begin{equation}
% S_n = \frac{a_1(1-q^n)}{1-q}
% \end{equation}
% \end{latexExample}

\subsection{Símbolos Especiais}
\subsubsection{Letras Gregas}
\begin{codeBox}
\begin{longtable}{ll|ll}
\verb|\alpha| & $\alpha$ & \verb|\beta| & $\beta$ \\
\verb|\gamma| & $\gamma$ & \verb|\Gamma| & $\Gamma$ \\
\verb|\delta| & $\delta$ & \verb|\Delta| & $\Delta$ \\
\verb|\epsilon| & $\epsilon$ & \verb|\varepsilon| & $\varepsilon$ \\
\verb|\zeta| & $\zeta$ & \verb|\eta| & $\eta$ \\
\verb|\theta| & $\theta$ & \verb|\vartheta| & $\vartheta$\\
\verb|\Theta| & $\Theta$ & \verb|\iota| & $\iota$ \\
\verb|\kappa| & $\kappa$ & \verb|\lambda| & $\lambda$ \\
\verb|\Lambda| & $\Lambda$ & \verb|\mu| & $\mu$ \\
\verb|\nu| & $\nu$ & \verb|\xi| & $\xi$ \\
\verb|\Xi| & $\Xi$ & \verb|\pi| & $\pi$ \\
\verb|\Pi| & $\Pi$ & \verb|\rho| & $\rho$ \\
\verb|\varrho| & $\varrho$ & \verb|\sigma| & $\sigma$ \\
\verb|\Sigma| & $\Sigma$ & \verb|\tau| & $\tau$ \\
\verb|\upsilon| & $\upsilon$ & \verb|\Upsilon| & $\Upsilon$ \\
\verb|\phi| & $\phi$ & \verb|\varphi| & $\varphi$ \\
\verb|\Phi| & $\Phi$ & \verb|\chi| & $\chi$ \\
\verb|\psi| & $\psi$ & \verb|\Psi| & $\Psi$ \\
\verb|\omega| & $\omega$ & \verb|\Omega| & $\Omega$ \\
\end{longtable}
\end{codeBox}

\subsubsection{Outro Símbolos Matemáticos}
\begin{codeBox}
\begin{longtable}{ll|l}
F.E.M (Força Eletro-Motriz) & \verb|\emf| & $\emf$ \\
Volume (V cortado) & \verb|\vol| & $\vol$ \\
Volume de Controle & \verb|\vc| & $\vc$ \\
Menor ou igual & \verb|\le| & $\le$ \\
Maior ou igual & \verb|\ge| & $\ge$ \\
União de conjunto & \verb|\cup| & $\cup$ \\
Disjunção de conjunto & \verb|\cap| & $\cap$ \\
Mais ou menos & \verb|\pm| & $\pm$ \\
Menos ou mais & \verb|\mp| & $\mp$ \\
Cunha (\emph{wedge}), e & \verb|\wedge| & $\wedge$ \\
Ou & \verb|\vee| & $\vee$ \\
Pertence & \verb|\in| & $\in$ \\
Não pertence & \verb|\notin| & $\notin$ \\
Contido em & \verb|\subset| & $\subset$ \\
Não contido em & \verb|\nsubset| & $\nsubset$ \\
\end{longtable}
\end{codeBox}

\subsection{Matrizes}

Matrizes podem ser inseridas usando o ambiente \verb|matrix| e suas variantes: \verb|p| para parênteses, \verb|b| para colchetes (do inglês \emph{brackets}), \verb|v| para barras verticais (útil para determinantes) e \verb|V| para barras verticais duplas.

\begin{latexExampleSide}
\[ \begin{matrix} 
a & b \\
c & d 
\end{matrix} \]
\end{latexExampleSide}

\begin{latexExampleSide}
\[ \begin{pmatrix} 
a & b \\
c & d 
\end{pmatrix} \]
\end{latexExampleSide}

\begin{latexExampleSide}
\[ \begin{bmatrix} 
a & b \\
c & d 
\end{bmatrix} \]
\end{latexExampleSide}

\begin{latexExampleSide}
\[ \begin{vmatrix} 
a & b \\
c & d 
\end{vmatrix} \]
\end{latexExampleSide}

\begin{latexExampleSide}
\[ \begin{Vmatrix} 
a & b \\
c & d 
\end{Vmatrix} \]
\end{latexExampleSide}

O alinhamento dos elementos pode ser feito utilizando os ambientes de matriz terminados em \verb|*| (asterisco) seguidos de um especificador de alinhamento: \verb|[r]| para direita e \verb|[l]| para esquerda.

\begin{latexExampleSide}
\[ \begin{bmatrix*}[r]
-7 & 0 \\
0 & 5 
\end{bmatrix*} \]
\end{latexExampleSide}

O alinhamento de números complexos é um pouco mais, bem... complexo. Utilizamos o comando \verb|\phantom| para colocar a unidade imaginária de forma ``invisível''.

\begin{latexExampleSide}
\[ \begin{bmatrix*}[r]
-7i & 0i \\
0\phantom{i} & 5\phantom{i} 
\end{bmatrix*} \]


\end{latexExampleSide}

\begin{latexExample}
This is a \LaTeX\ example:
\begin{equation}
\sum\limits_{i=1}^n i = \frac{n(n+1)}{2}.
\end{equation}
\end{latexExample}

\end{document}
